
\chapter{परोपकार में हमारा ही उपकार है}

यह विचार करने के पहले कि कर्तव्यनिष्ठा हमें आध्यात्मिक उन्नति में किस प्रकार सहायता पहुँचाती है, मैं आप लोगों को संक्षेप में यह भी बता देना चाहता हूँ कि भारतवर्ष में जिसे हम कर्म कहते हैं, उसका एक दूसरा पहलू और कौनसा है~। प्रत्येक धर्म के तीन विभाग होते हैं~। प्रथम दार्शनिक, दूसरा पौराणिक और तीसरा कर्मकाण्ड~। दार्शनिक भाग तो असल में प्रत्येक धर्म का सार है~। पौराणिक भाग इस दार्शनिक भाग की व्याख्यास्वरूप है; उसमें महापुरुषों की कम या अधिक काल्पनिक जीवनी तथा अलौकिक विषयसम्बन्धी कथाओं एवं आख्यायिकाओं द्वारा इसी दर्शन को उत्तम रूप से समझाने की चेष्टा की गयी है~। और कर्मकाण्ड इस दर्शन का ही और भी स्थूल रूप है, जिससे वह सर्वसाधारण की समझ में आ सके~। वास्तव में अनुष्ठान दर्शन का ही एक स्थूलतर रूप है~। यह अनुष्ठान ही कर्म है~। प्रत्येक धर्म में इसकी आवश्यकता है, क्योंकि जब तक हम आध्यात्मिक जीवन में बहुत उन्नत न हो जायँ, तब तक सूक्ष्म आध्यात्मिक तत्त्व को समझ ही नहीं सकते~। मनुष्य को अपने तई यह मान लेना सरल है कि वह कोई भी बात समझ सकता है~। परन्तु जब वह उसे अमल में लाने की चेष्टा करता है, तो उसे मालूम होता है कि सूक्ष्म भावों को ठीकठीक समझना तथा उन्हें हृदयंगम करना बड़ा ही कठिन है~। इसीलिए तो प्रतीक विशेष रूप से सहाय्यक होते हैं~। प्रतीक द्वारा सूक्ष्म विषयों को समझने की जो प्रणाली है, उसे हम किसी प्रकार त्याग नहीं सकते~। अत्यन्त प्राचीन समय से ही प्रतीकों का प्रयोग प्रत्येक धर्म में होता रहा है~। एक दृष्टि से यदि कहा जाय कि हम प्रतीक के बिना किसी बात को सोच ही नहीं सकते, तो ठीक ही है~। शब्द भी तो हमारे विचार के प्रतीक ही हैं~। दूसरी दृष्टि से कहा जा सकता है कि संसार की प्रत्येक चीज प्रतीक के रूप में देखी जा सकती है~। सारा संसार ही प्रतीक है और उसके पीछे मूलतत्त्व रूप में ईश्वर विराजमान हैं~। इस प्रकार का प्रतीक केवल मनुष्य द्वारा उत्पन्न किया हुआ ही नहीं है~। और न ऐसी ही बात है कि एक धर्म के कुछ अनुयायी एक साथ बैठ गये और बस उन्होंने कुछ प्रतीकों की कल्पना कर डाली~। धर्म के प्रतीक की उत्पत्ति स्वाभाविक रूप से होती है~। नहीं तो ऐसा क्यों है कि प्रायः सभी मनुष्यों के मन में कुछ विशेष प्रतीक कुछ विशिष्ट भावों से सदा सम्बद्ध रहते है? कुछ प्रतीक तो सभी जगह पाये जाते हैं~। तुम में से अनेकों की यह धारणा है कि क्रास का चिह्न सर्वप्रथम ईसाई धर्म के साथ प्रचलित हुआ, परन्तु वास्तव में तो वह ईसाई धर्म के बहुत पहले से, मूसा के भी जन्म के पहले, वेदों के आविर्भाव के भी पहले यहाँ तक कि मानवी कार्यकलापों का किसी प्रकार का इतिहास लिपिबद्ध होने के भी पहले से विद्यमान है~। ऐझ्टिक्स तथा फिनिशिन्स जातियों में भी इसकी मौजूदगी का प्रमाण मिलता है; प्रायः प्रत्येक जाति में इसका अस्तित्व था। इतना ही नहीं बल्कि ऐसा भी प्रतीत होता है कि क्रॉस पर लटके हुए महापुरुष का प्रतीक लगभग प्रत्येक जाति में प्रचलित था~। इसी प्रकार सारे संसार में वृत्त भी एक उत्कृष्ट प्रतीक माना गया है~। फिर इसके अतिरिक्त सारे संसार में सबसे अधिक प्रचलित स्वस्तिक (\includegraphics[scale=0.8]{"images/1.png"}) भी एक प्रतीक रहा है~। एक समय ऐसी धारणा थी कि बौद्ध इसे अपने साथ-साथ सारे संसार भर में ले गये; परन्तु पता चलता है कि बौद्ध धर्म के सदियों पहले कई जातियों में इसका प्रचार था~। प्राचीन बैबिलोन तथा मिश्र देश में भी यह पाया जाता था~। इस सबसे क्या प्रकट होता है? यही कि ये सब प्रतीक केवल काल्पनिक या स्वेच्छाकृत ही नहीं थे~। इनका कोई न कोई विशेष कारण अवश्य रहा होगा, उनमें तथा मानवी मन में कोई स्वाभाविक सम्बन्ध रहा होगा~। इसी प्रकार भाषा भी कोई कृत्रिम वस्तु नहीं है~। ऐसी बात नहीं कि कुछ लोगों ने एक साथ बैठकर यह तय कर लिया कि कुछ विशेष भाव कुछ विशेष शब्दों द्वारा प्रकट किये जायें और बस भाषा की उत्पत्ति हो गयी~। नहीं, भाषा की उत्पत्ति इस प्रकार नहीं हुई~। कोई भी भाव अपने आनुषंगिक शब्द बिना नहीं रह सकता, और न कोई शब्द ही अपने आनुषंगिक भाव बिना रह सकता है। शब्द और भाव स्वभावतः अविच्छेद्य हैं~। भावों को प्रकट करने के लिए शब्दप्रतीक अथवा वर्णप्रतीक - किसी का भी व्यवहार किया जा सकता है~। गूँगों और बहरों को शब्दप्रतीक किसी प्रकार की सहायता नहीं पहुँचा सकता, उन्हें किसी दूसरे प्रतीक की सहायता लेनी पड़ती है~। मन में उठनेवाले प्रत्येक विचार का एक दूसरा अंश भी होता है, और वह है - आकृति~। इसे संस्कृत दर्शन में ‘नामरूप’ कहते हैं~। जिस प्रकार कृत्रिम उपाय द्वारा एक भाषा नहीं उत्पन्न की जा सकती, उसी प्रकार कृत्रिम उपायों से प्रतीक का निर्माण भी नहीं किया जा सकता~। संसार में कर्मकाण्ड के सहकारी जो-जो प्रतीक हैं, वे मानवजाति के धार्मिक विचारों के एक एक बाह्य रूप हैं~। यह कह देना बहुत सरल है कि अनुष्ठानों, मन्दिरों तथा अन्य बाह्य आडम्बरों की कोई आवश्यकता नहीं, और यह बात तो कल का एक छोकरा भी कहा करता है~। परन्तु सरलतापूर्वक यह कोई भी देख सकता है कि जो लोग मन्दिर में जाकर पूजा करते हैं वे उन लोगों की अपेक्षा, जो ऐसा नहीं करते, कई बातों में कही भिन्न होते हैं~। भिन्न-भिन्न धर्मों के साथ जो विशिष्ट मन्दिर, अनष्ठान और अन्य स्थूल क्रियाकलाप जड़े हुए हैं, वे उन-उन धर्मावलम्बियों के मन में उन सब भावों को जागृत कर देते हैं, जिनके कि ये मन्दिर-अनुष्ठानादि स्थूल प्रतीक स्वरूप हैं~। अतएव अनुष्ठानों एवं प्रतीकों को एकदम उड़ा देना उचित नहीं~। इन सब विषयों का अध्ययन एवं अभ्यास स्वभावतः कर्मयोग का ही एक अंग है~।

इस कर्मयोग के और भी कई पहलू हैं~। इनमें से एक है - ‘भाव’ तथा ‘शब्द’ के सम्बन्ध को जानना एवं यह भी ज्ञान प्राप्त करना कि शब्दशक्ति से क्या-क्या हो सकता है~। प्रत्येक धर्म शब्द की शक्ति को मानता है; यहाँ तक कि किसी-किसी धर्म की तो यह धारणा है कि समस्त सृष्टि ‘शब्द’ से ही निकली है~। ईश्वर के संकल्प का बाह्य आकार ‘शब्द’ है और चूँकि ईश्वर ने सृष्टिरचना के पूर्व संकल्प एवं इच्छा की थी, इसलिए सृष्टि ‘शब्द’ से ही निकली है~। इस जड़वादमय जीवन के कोलाहल में हमारे स्नायुओं में भी जड़ता आ गयी है~। ज्यों-ज्यों हम बूढ़े होते जाते हैं और संसार की ठोकरें खाते जाते हैं, त्यों-त्यों हममें अधिकाधिक जड़ता आती जाती है और फलस्वरूप, हमारे चारों ओर निरन्तर हमारे ध्यान को आकर्षित करनेवाली जो सारी घटनाएँ होती रहती हैं, उनके प्रति हम उदासीन रहकर उनसे कोई शिक्षा ग्रहण नहीं कर पाते~। परन्तु कभी-कभी ऐसा भी अवश्य होता है कि मनुष्य का स्वाभाविक भाव प्रबल हो उठता है और हम इन साधारण सांसारिक घटनाओं का रहस्य जानने का यत्न करने लगते हैं तथा उन पर आश्चर्यचकित हो जाते हैं~। इस प्रकार आश्चर्यचकित होना ही ज्ञानलाभ की पहली सीढ़ी है~।

शब्द के उच्च दार्शनिक तथा धार्मिक महत्त्व को छोड़ देने पर भी हम देखते हैं कि हमारे इस जीवननाटक में शब्दप्रतीक का विशेष स्थान है~। मैं तुमसे बातचीत कर रहा हूँ, तुम्हें स्पर्श नहीं कर रहा हूँ~। पर मेरे शब्द द्वारा उत्पन्न वायु के स्पन्दन तुम्हारे कान में जाकर तुम्हारे कर्मस्नायुओं को स्पर्श करते हैं और उससे तुम्हारे मन में असर पैदा होता है~। इसे तुम रोक नहीं सकते~। भला सोचो तो, इससे अधिक आश्चर्यजनक बात और क्या हो सकती है? एक मनुष्य दूसरे को ‘बेवकूफ’ कह देता है और बस इतने से ही वह दूसरा मनुष्य उठ खड़ा होता है और अपनी मुठ्ठी बाँधकर उसकी नाक पर एक घूँसा जमा देता है~। देखो तो शब्द में कितनी शक्ति है! एक स्त्री बिलखबिलखकर रो रही है; इतने में एक दूसरी स्त्री आ जाती है और वह उससे कुछ सान्त्वना के शब्द कहती है~। प्रभाव यह होता है कि वह रोती हुई स्त्री उठ बैठती है, उसका दुःख दूर हो जाता है और वह मुसकराने भी लगती है~। देखो तो शब्द में कितनी शक्ति है! उच्च दर्शन में जिस प्रकार शब्दशक्ति का परिचय मिलता है, उसी प्रकार साधारण जीवन में भी~। इस शक्ति के सम्बन्ध में विशेष विचार और अनुसन्धान न करते हुए भी हम रात-दिन इस शक्ति को उपयोग में ला रहे हैं~। इस शक्ति के स्वरूप को जानना तथा इसका उत्तम रूप से उपयोग करना भी कर्मयोग का एक अंग है~।

दूसरों के प्रति हमारे कर्तव्य का अर्थ है - दूसरों की सहायता करना, संसार का भला करना~। अब प्रश्न उठता है कि हम संसार का भला क्यों करे? वास्तव में बात यह है कि ऊपर से तो हम संसार का उपकार करते हैं, परन्तु असल में हम अपना ही उपकार करते हैं~। हमें सदैव संसार का उपकार करने की चेष्टा करनी चाहिए, और कार्य करने में यही हमारा सर्वोच्च उद्देश्य होना चाहिए~। परन्तु यदि ध्यानपूर्वक देखा जाय, तो प्रतीत होगा कि इस संसार को हमारी सहायता की बिलकुल आवश्यकता नहीं~। यह संसार इसलिए नहीं बना कि हम अथवा तुम आकर इसकी सहायता करें~। एक बार मैंने एक उपदेश पढ़ा था, वह इस प्रकार था - “यह सुन्दर संसार बड़ा अच्छा है, क्योंकि इसमें हमें दूसरों की सहायता करने के लिए समय तथा अवसर मिलता है~।” ऊपर से तो यह भाव सचमुच सुन्दर है, परन्तु यह कहना कि संसार को हमारी सहायता की आवश्यकता है, क्या घोर ईश्वर-निन्दा नहीं है? यह सच है कि संसार में दुःखकष्ट बहुत है, और इसलिए लोगों की सहायता करना हमारे लिए सर्वश्रेष्ठ कार्य है; परन्तु आगे चलकर हम देखेंगे कि दूसरों की सहायता करने का अर्थ है अपनी ही सहायता करना~। मुझे स्मरण है, एक बार जब मैं छोटा था, तो मेरे पास कुछ सफेद चूहे थे~। वे चूहे एक छोटेसे सन्दूक में रखे गये थे और उस सन्दूक के भीतर उनके लिये छोटे-छोटे चक्के बना दिये गये थे~। जब चूहे उन चक्कों को पार करना चाहते, तो वे चक्के वहीं के वहीं घूमते रहते, और वे बेचारे चूहे कभी भी बाहर नहीं निकल पाते~। बस यही हाल संसार का तथा संसार के प्रति हमारी सहायता का है~। उपकार केवल इतना ही होता है कि हमें नैतिक शिक्षा मिलती है~। संसार न तो अच्छा है, न बुरा बात इतनी ही है कि प्रत्येक मनुष्य अपने लिए अपना-अपना संसार बना लेता है~। यदि एक अन्धा संसार के बारे में सोचने लगे, तो वह उसके समक्ष या तो मुलायम या कड़ा प्रतीत होगा, अथवा शीत या उष्ण~। हम सुख या दुःख की समष्टि मात्र हैं; यह हमने अपने जीवन में सैकड़ों बार अनुभव किया है~। बहुधा नौजवान आशावादी होते हैं, और वृद्ध निराशावादी~। तरुण के सामने अभी उसका सारा जीवन पड़ा है~। परन्तु वृद्ध की केवल यही शिकायत रहती है कि उसका समय निकल गया; कितनी ही अपूर्ण इच्छाएँ उसके हृदय में मचलती रहती हैं, जिन्हें पूर्ण करने की शक्ति उसमें आज नहीं~। परन्तु हैं दोनों ही मूर्ख~। हमारी मानसिक अवस्था के अनुसार ही हमें यह संसार भला या बुरा प्रतीत होता है~। स्वयं यह न तो भला है, न बुरा~। अग्नि स्वयं न अच्छी है, न बुरी~। जब यह हमें गरम रखती है, तो हम कहते हैं, ‘यह कितनी सुन्दर है!’ परन्तु जब इससे हमारी उँगली जल जाती है, तो इसे हम दोष देते हैं~। परन्तु फिर भी स्वयं न तो यह अच्छी है, न बुरी~। जैसा हम इसका उपयोग करते हैं, तदनुरूप यह अच्छी या बुरी बन जाती है~। यही हाल इस संसार का भी है~। संसार स्वयं पूर्ण है~। पूर्ण होने का अर्थ यह है कि उसमें अपने सब प्रयोजनों को पूर्ण करने की क्षमता है~। हमें यह निश्चित जान लेना चाहिए कि हमारे बिना भी यह संसार बड़े मजे से चलता जाएगा; हमें उसकी सहायता करने के लिए माथापच्ची करने की आवश्यकता नहीं~।

परन्तु फिर भी हमें सदैव परोपकार करते ही रहना चाहिए~। यदि हम सदैव यह ध्यान रखें कि दूसरों की सहायता करना एक सौभाग्य है, तो परोपकार करने की इच्छा एक सर्वोत्तम प्रेरणाशक्ति है~। एक दाता के ऊँचे आसन पर खड़े होकर और अपने हाथ में दो पैसे लेकर यह मत कहो, “ऐ भिखारी, ले, यह मैं तुझे देता हूँ~।” परन्तु तुम स्वयं इस बात के लिए कृतज्ञ होओ कि तुम्हें वह निर्धन व्यक्ति मिला, जिसे दान देकर तुमने स्वयं अपना उपकार किया~। धन्य पानेवाला नहीं होता, देनेवाला होता है~। इस बात के लिए कृतज्ञ होओ कि इस संसार में तुम्हें अपनी दयालुता का प्रयोग करने और इस प्रकार पवित्र एवं पूर्ण होने का अवसर प्राप्त हुआ~। समस्त भले कार्य हमें शुद्ध बनने तथा पूर्ण होने में सहायता करते हैं~। और सच पूछो तो हम अधिक से अधिक कर ही कितना सकते हैं? या तो एक अस्पताल बनवा देते हैं, सड़के बनवा देते हैं या सदावर्त खुलवा देते हैं, बस इतना ही तो? हम गरीबों के लिए एक कोष खोल देते हैं, दस-बीस लाख रुपये इकट्ठा कर लेते हैं~। उसमें से पाँच लाख का एक अस्पताल बनवा देते हैं, पाँच लाख नाचतमाशे में फूँक देते हैं और शेष का आधा कर्मचारी लूट लेते हैं; बाकी जो बचता है, वह किसी तरह गरीबों तक पहुँचता है! परन्तु उतने से हुआ क्या? आँधी का एक झोंका तो तुम्हारी इन सारी इमारतों को पाँच मिनट में नष्ट कर दे सकता है - फिर तुम क्या करोगे? एक भूकम्प तो तुम्हारी तमाम सड़कों, अस्पतालों, नगरों और इमारतों को धूल में मिला दे सकता है~। अतएव इस प्रकार की संसार की सहायता करने की खोखली बातों को हमें छोड़ देना चाहिए~। यह संसार न तो तुम्हारी सहायता का भूखा है और न मेरी~। परन्तु फिर भी हमें निरन्तर कार्य करते रहना चाहिए, निरन्तर परोपकार करते रहना चाहिए~। क्यों? - इसलिए कि इससे हमारा ही भला है~। यही एक साधन है, जिससे हम पूर्ण बन सकते हैं~। यदि हमने किसी भिखारी को कुछ दिया हो, तो वास्तव में उसके ऊपर हमारा कुछ नहीं आता, हमारे ऊपर ही उसका आता है, हम पर उसका आभार है, क्योंकि उसने हमें इस बात का अवसर दिया कि हम अपनी दया उस पर काम में ला सकें~। यह सोचना निरी भूल है कि हमने संसार का भला किया, अथवा कर सकते हैं, या यह कि हमने अमुक-अमुक व्यक्तियों की सहायता की~। यह निरी मूर्खता का विचार है; और मूर्खता के विचारों से दुःख उत्पन्न होता है~। हम कभी-कभी सोचते हैं कि हमने अमुक मनुष्य की सहायता की और इसलिए आशा करते हैं कि वह हमें धन्यवाद दे; पर जब वह हमें धन्यवाद नहीं देता, तो उससे हमें दुःख होता है~। हम जो कुछ करें, उसके बदले में किसी भी बात की आशा क्यों रखें? बल्कि उलटे हमें उसी मनुष्य के प्रति कृतज्ञ होना चाहिए, जिसकी हम सहायता करते हैं - उसे साक्षात् नारायण मानना चाहिए~। मनुष्य की सहायता द्वारा ईश्वर की उपासना करना क्या हमारा परम सौभाग्य नहीं है? यदि हम वास्तव में अनासक्त हैं, तो हमें यह वृथा प्रत्याशाजनित कष्ट क्यों होना चाहिए? अनासक्त होने पर तो हम प्रसन्नतापूर्वक संसार में भलाई कर सकते हैं~। अनासक्ति से किये हुए कार्य से कभी भी दुःख अथवा अशान्ति नहीं आएगी~। वैसे तो संसार में अनन्त काल तक सुखदुःख का चक्र चलता ही रहेगा, फिर हम इसकी सहायता के लिए कुछ करें या न करें, उससे कुछ बनने बिगड़ने का नहीं~। दृष्टान्तस्वरूप एक कहानी सुनो:-

एक गरीब आदमी को कुछ रुपये की जरूरत पड़ी~। उसे कहीं से यह मालूम हो गया कि यदि वह किसी भूत को अपने वश में कर ले, तो वह उससे जो चाहे मँगवा सकता है~। निदान उसे एक भूत ढूँढ़ने की सूझी~। वह किसी ऐसे आदमी को ढूँढ़ने लगा, जिससे उसे एक भूत मिल जाय~। ढूँढ़ते-ढूँढ़ते उसे एक साधु मिले~। इन साधु के पास बड़ी शक्तियाँ थीं और उसने उनसे सहायता की याचना की~। साधु ने उससे पूछा, “तुम भूत का क्या करोगे?” उसने उत्तर दिया, “महाराज मैं भूत इसलिए चाहता हूँ कि वह मेरा काम कर दे~। कृपा कर मुझे उसको प्राप्त करने का ढंग बता दीजिये~। मुझे उसकी बड़ी जरूरत है~।” साधु बोले, “देखो, तुम इस झमेले में मत पड़ो, अपने घर लौट जाओ~।” दूसरे दिन वह आदमी साधु के पास फिर गया और बहुत रोने-गाने लगा~। उसने कहा, “महाराज, मुझे एक भूत दे ही दीजिये न~। मुझे बड़ी आवश्यकता है~।” अन्त में साधु कुछ चिढ़से गये और उन्होंने कहा, “अच्छा, लो, यह मन्त्र लो, इसका जप करने से एक भूत प्रकट होगा और फिर उससे तुम जो काम कहोगे, वही करेगा; परन्तु देखो, होशियार रहना~। ये बड़े भयंकर प्राणी होते है~। उसे निरन्तर काम में लगाये रखना~। यदि कभी वह खाली रहा, तो तुम्हारी जान ही ले लेगा~।” तब उस मनुष्य ने कहा, “यह कौन कठिन बात है? मैं तो उसे इतना काम दे दूँ कि उसके जीवन भर खत्म न हो~।” इसके बाद वह आदमी एक वन में चला गया और मन्त्र का जप करने लगा~। कुछ देर तक जप करने के बाद उसके सामने विकराल दाँतोंवाला एक बड़ा भयंकर भूत प्रकट हुआ~। भूत ने कहा, “देखो, मैं भूत हूँ~। तुम्हारे मन्त्र ने मुझे जीत लिया है~। परन्तु देखो तुम्हें मुझको निरन्तर काम में लगाये रहना होगा, क्योंकि ज्योंही मुझे थोड़ासा भी अवकाश मिला कि मैं तुम्हारी जान ले लूँगा~।” आदमी बोला, “ठीक है, जाओ, मेरे लिए एक महल तैयार करो~।” भूत ने जवाब दिया, “लो, हो गया महल तैयार है~।” आदमी ने कहा, “जाओ, मेरे लिए धन ले आओ~।” भूत बोला, “लो, धन भी तैयार है~।” फिर आदमी ने कहा, “यह जंगल काट डालो और यहाँ एक शहर बसा दो~।” भूत बोला, “लो, यह भी हो गया~। अब और क्या करूँ बतलाओ?” अब तो यह आदमी बड़ा घबड़ाने लगा; उसने मन में सोचा, “अब तो मेरे पास कोई काम नहीं है, जो मैं इसे करने को कहूँ~। यह तो प्रत्येक काम क्षण भर में ही कर डालता है~।” भूत इधर गरजकर बोला, “देखो, मुझे जल्दी कुछ काम करने को दो, नहीं तो मैं तुम्हें खा जाऊँगा~।” बेचारा आदमी अब कोई काम सोच ही न सका और मारे भय के थर-थर काँपने लगा~। अब तो वह बेतहाशा भागा और भागते भागते उन्हीं साधु के पास पहुँचा और वहाँ जाकर गिड़गिड़ाने लगा, “महाराज, रक्षा कीजिये, रक्षा कीजिये, मेरी जान बचाइये~।” साधु ने पूछा, “कहो, क्या हुआ”; मनुष्य ने उत्तर दिया, “अब मैं क्या करूँ? अब तो मेरे पास उस भूत को देने के लिए कोई भी काम शेष नहीं है~। मैं उससे जो कुछ भी करने को कहता हूँ, वह क्षण भर में कर डालता है, और जब उसके पास कोई काम नहीं रह जाता, तो मुझे खा डालने की धमकी देता है~।” इतने में ही वह भूत वहाँ आ पहुँचा, और कहने लगा, “अब तो मैं तुम्हें खा जाऊँगा~।” और सचमुच वह उसे खा जाता! आदमी मारे डर के काँपने लगा और उसने साधु से अपने प्राणों की भिक्षा मांगी~। साधु ने कहा, “अच्छा, मैं तुम्हें रास्ता बताता हूँ~। देखो, उस कुत्ते की पूँछ टेढ़ी है, अपनी तलवार निकालो और यह पूँछ काटकर इस भूत को दे दो और उससे कहो कि इसे सीधी कर दे~।” आदमी ने झट से कुत्ते की पूँछ काट ली और उसे भूत को देकर कहा, “लो, इसे सीधी करके मुझे दो~।” भूत ने पूँछ ले ली और उसे बड़ी सावधानी से सीधी की, पर ज्यों ही उसने उसको सीधी करके छोड़ दिया, त्यों ही वह फिर से टेढ़ी हो गयी~। भूत ने दुबारा कोशिश की परन्तु ज्योंही उसने छोड़ दी, त्यों ही वह फिर टेढ़ी हो गयी~। उसने तीसरी बार फिर प्रयत्न किया, परन्तु वह फिर टेढ़ी की टेढ़ी हो गयी~। इस प्रकार वह कई दिनों तक प्रयत्न करता रहा, यहाँ तक कि वह थक गया और बोला, “मुझे ऐसा कष्ट तो अपने जीवन भर में कभी नहीं हुआ~। मैं एक बड़ा पुराना भूत हूँ, ऐसी मुसीबत में मैं कभी नहीं पड़ा~।” अब तो वह भूत उस आदमी से कहने लगा, “आओ, भाई, हम-तुम समझौता कर लें~। तुम मुझे छोड़ दो, और मैंने अब तक तुम्हें जो कुछ दिया है, वह सब अपने पास ही रखे रहो~। मैं वादा करता हूँ, अब आगे तुम्हें किसी प्रकार का कष्ट न दूँगा~।” यह सुन वह आदमी बड़ा खुश हुआ और बड़ी प्रसन्नतापूर्वक उसने इस समझौते को स्वीकार कर लिया~।

हमारा यह संसार भी बस कुत्ते की टेढ़ी पूँछ के ही समान है; सैकड़ों वर्ष से लोग इसे सीधा करने का प्रयत्न कर रहे हैं, परन्तु ज्यों ही वे इसे छोड़ देते हैं, त्यों ही यह फिर से टेढ़ा का टेढ़ा हो जाता है~। इसके अतिरिक्त और क्या हो सकता है? मनुष्य पहले यह जान ले कि आसक्तिरहित होकर उसे किस प्रकार कर्म करना चाहिए, तभी वह दुराग्रह और मतान्धता से परे हो सकता है~। जब हमें यह ज्ञान हो जाएगा कि संसार कुत्ते की टेढ़ी दुम की तरह है और कभी भी सीधा नहीं हो सकता, तब हम दुराग्रही नहीं होगें~। यदि संसार में यह दुराग्रह, यह कट्टरता न होती, तो अब तक यह बहुत उन्नति कर लेता~। यह सोचना भूल है कि धर्मान्धता द्वारा मानवजाति की उन्नति हो सकती है~। बल्कि उल्टे, यह तो हमें पीछे हटानेवाली शक्ति है, जिससे घृणा और क्रोध उत्पन्न होकर मनुष्य एक दूसरे से लड़नेभिड़ने लगते हैं और सहानुभूतिशून्य हो जाते हैं~। हम सोचते हैं कि जो कुछ हमारे पास है अथवा जो कुछ हम करते हैं, वही संसार में सर्वश्रेष्ठ है, और जो कुछ हम नहीं करते अथवा जो कुछ हमारे पास नहीं है, वह एक कौड़ी मूल्य का भी नहीं~। अतएव जब कभी तुममें दुराग्रह का यह भाव आये, तो सदैव कुत्ते की टेढ़ी पूँछ का दृष्टान्त स्मरण कर लिया करो~। तुम्हें अपने-आपको संसार के बारे में चिन्तित बना लेने की कोई आवश्यकता नहीं - तुम्हारी सहायता के बिना भी यह चलता ही रहेगा~। जब तुम दुराग्रह और मतान्धता से परे हो जाओगे, तभी तुम अच्छी तरह कार्य कर सकोगे~। जो ठण्डे मस्तिष्कवाला और शान्त है, जो उत्तम ढंग से विचार करके कार्य करता है, जिसके स्नायु सहज ही उत्तेजित नहीं हो उठते तथा जो अत्यन्त प्रेम और सहानुभूति-सम्पन्न है, केवल वही व्यक्ति संसार में महान कार्य कर सकता है और इस तरह उससे अपना भी कल्याण कर सकता है~। दुराग्रही व्यक्ति मूर्ख और सहानुभूतिशून्य होता है~। वह न तो कभी संसार को सीधा कर सकता है और न स्वयं ही शुद्ध एवं पूर्ण हो सकता है~।

आज के व्याख्यान का सारांश यह है~। सर्वप्रथम हमें यह ध्यान में रखना चाहिए कि हम ही संसार के ऋणी हैं, संसार हमारा ऋणी नहीं~। यह तो हमारा सौभाग्य है कि हमें संसार में कुछ कार्य करने का अवसर मिलता है~। संसार की सहायता करने से हम वास्तव में स्वयं का ही कल्याण करते हैं~। दूसरी बात यह है कि इस विश्व के एक ईश्वर हैं~। यह बात सच नहीं कि यह संसार पिछड़ रहा है और इसे तुम्हारी अथवा मेरी सहायता की आवश्यकता है~। ईश्वर सर्वत्र विराजमान हैं~। वे अविनाशी, सतत क्रियाशील और जाग्रत हैं~। जब सारा विश्व सोता है, तब भी वे जागते रहते हैं~। वे निरन्तर कार्य में लगे हुए हैं~। संसार के समस्त परिवर्तन और विकास उन्ही के कार्य हैं~। तीसरी बात यह है कि हमें किसी से घृणा नहीं करनी चाहिए~। यह संसार सदैव ही अच्छे और बुरे का मिश्रणस्वरूप रहेगा~। हमारा कर्तव्य है कि हम दुर्बल के प्रति सहानुभूति रखें और एक अन्यायी के प्रति भी प्रेम रखें~। यह संसार तो चरित्रगठन के लिए एक विशाल व्यायामशाला है~। इसमें हम सभी को अभ्यासरूप कसरत करनी पड़ती है, जिससे हम आध्यात्मिक बल से अधिकाधिक बलवान बनते रहें~। चौथी बात यह है कि हममें किसी प्रकार का भी दुराग्रह नहीं होना चाहिए, क्योंकि दुराग्रह प्रेम का विरोधी है~। बहुधा दुराग्रहियों को तुम यह गाल बजाते सुनोगे, “हमें पापी से घृणा नहीं है, हमें तो घृणा पाप से है~।” परन्तु यदि कोई मुझे एक ऐसा मनुष्य दिखा दे जो सचमुच पाप और पापी में भेद कर सकता है, तो ऐसे मनुष्य को देखने के लिए मैं कितनी भी दूर जाने को तैयार हूँ~। मुख से ऐसा कहना सरल है~। यदि हम द्रव्य और उसके गुण में भलीभाँति भेद कर सकें, तो हम पूर्ण हो जायँ~। पर इसे अमल में लाना इतना सरल नहीं~। हम जितने ही शान्तचित्त होंगे और हमारे स्नायु जितने शान्त रहेंगे हम उतने ही अधिक प्रेमसम्पन्न होंगे और हमारा कार्य भी उतना ही अधिक श्रेष्ठ होगा~।

